\chapter{习题参考解答}
\label{app:solutions}
这些解答给出判断依据和关键计算。凡是通过变换求解的题目，最后都回到原约束检查结果。

\section*{题~\ref{ex:units}：单位与带宽}
接收信号功率为 $gp=10^{-9}\,\mathrm W$，带内噪声为 $N_0B=10^{-11}\,\mathrm W$，故信噪比为 $100$，即 $20\,\mathrm{dB}$。速率是
\[
 R=B\log_2\!\left(1+\frac{gp}{N_0B}\right)
 =10^6\log_2(101)\approx6.6582\times10^6\,\mathrm{bit/s}.
\]
带宽加倍后信噪比降为 $50$，速率为 $2\times10^6\log_2(51)\approx11.3449\times10^6\,\mathrm{bit/s}$。带宽增加既增加每秒自由度，也增加噪声；在固定总功率下速率并不按带宽成倍增加。自然对数计算的速率为 $\mathrm{nat/s}$，除以 $\ln2$ 才得到 $\mathrm{bit/s}$。

\section*{题~\ref{ex:nonconcave}：联合凹性的反例}
两个端点的总速率均为 $\ln2\approx0.6931$；中点的值为 $2\ln(5/4)\approx0.4463$，小于端点值的平均数，违反凹函数的不等式。固定他人功率后对自身功率的凹性，只描述一个坐标方向。联合凹性要求任意两个可行功率向量之间都满足凹性不等式，交叉干扰会破坏这一要求。因此最大化此总速率不构成凹目标最大化的凸优化问题。

\section*{题~\ref{ex:water}：注水与互补松弛}
最大化目标严格凹且各分量单调增加，故最优时用完功率。把问题写成负目标的最小化，乘子 $\lambda\geq0$ 对应总功率，$\nu_i\geq0$ 对应 $-p_i\leq0$，驻点条件为
\[
 -\frac1{\alpha_i+p_i}+\lambda-\nu_i=0,
 \qquad \nu_i p_i=0.
\]
正功率信道满足 $p_i=1/\lambda-\alpha_i$，故 $p_i=[\mu-\alpha_i]_+$，其中水位 $\mu=1/\lambda$。若前三条信道激活，$3\mu-(1+2+4)=10$，得到
\[
 \mu=\frac{17}{3},\qquad
 p^*=\left(\frac{14}{3},\frac{11}{3},\frac{5}{3},0\right),
 \qquad\lambda=\frac3{17}.
\]
水位位于 $4$ 和 $8$ 之间，激活集合自洽。第四条信道的乘子为 $\nu_4=3/17-1/8=7/136>0$，满足 KKT 条件。严格凹性保证上述可行 KKT 解唯一。

\section*{题~\ref{ex:power}：最小可行功率}
SINR 约束等价于 $p_1\geq0.4p_2+0.2$、$p_2\geq0.4p_1+0.2$。最小化 $p_1+p_2$，再加上非负及峰值约束，就是线性规划。这里
\[
 F=\begin{bmatrix}0&0.4\\0.4&0\end{bmatrix},\qquad
 u=\begin{bmatrix}0.2\\0.2\end{bmatrix},\qquad
 \rho(F)=0.4<1.
\]
无峰值约束时，分量最小可行功率为 $(I-F)^{-1}u=(1/3,1/3)^{\T}$。直接代入得 $(1/3)/(0.2/3+0.1)=2$。任何可行向量都不小于这个分量最小解，因此峰值为 $0.3$ 时不可行。

\section*{题~\ref{ex:soc}：相位选择与锥约束}
把 $w_k$ 乘以单位模复数 $\exp[-\mathrm i\arg(h_k^{\HH}w_k)]$，可使其期望信号为非负实数。该操作保持所有 $|h_\ell^{\HH}w_k|$、$\|w_k\|_2$ 不变；题设的目标和约束因此均不改变。固定相位后原约束等价于
\[
 \left\|\begin{bmatrix}(h_k^{\HH}w_j)_{j\ne k}\\\sigma_k\end{bmatrix}\right\|_2
 \leq\frac{\Rea(h_k^{\HH}w_k)}{\sqrt{\gamma_k}},
 \qquad \Ima(h_k^{\HH}w_k)=0.
\]
右侧由锥约束自动要求非负。复向量的范数可用其实部、虚部的堆叠表示成实二阶锥。多播时同一个 $w$ 同时出现在多个期望信号 $h_k^{\HH}w$ 中；一个整体相位通常无法让它们全部为非负实数，因此不能由此得到相同的等价凸化。

\section*{题~\ref{ex:robust}：最坏误差的方向}
由 Cauchy--Schwarz 不等式，$\delta^{\T}x\leq\varepsilon\|x\|_2$。当 $x\ne0$ 时，取 $\delta=\varepsilon x/\|x\|_2$ 即达到上界，故
\[
 \max_{\|\delta\|_2\leq\varepsilon}(\hat a+\delta)^{\T}x
 =\hat a^{\T}x+\varepsilon\|x\|_2.
\]
最坏误差沿着使左端增长最快的方向。这证明了鲁棒约束的等价形式。若 $x=0$，最大值为 $0$，任意允许误差均达到该值，约束退化为 $0\leq b$。

\section*{题~\ref{ex:prox}：复数软阈值}
若 $v\ne0$，固定 $|z|=r$ 后，$|z-v|^2=r^2+|v|^2-2\Rea(z^*v)$ 在 $z$ 与 $v$ 同相时最小。剩下的标量问题是
$\min_{r\geq0}\frac12(r-|v|)^2+\lambda r$，其解为 $r=[|v|-\lambda]_+$。于是
\[
 \prox_{\lambda|\cdot|}(v)=
 \begin{cases}
 (1-\lambda/|v|)v,&|v|>\lambda,\\
 0,&|v|\leq\lambda.
 \end{cases}
\]
给定数值的答案为 $\frac35(3+4\mathrm i)=1.8+2.4\mathrm i$。当 $v\in\R$ 时，它成为 $\operatorname{sign}(v)[|v|-\lambda]_+$；$v=0$ 时解为 $0$。

\section*{题~\ref{ex:dual}：对偶上升而非下降}
对偶函数 $d(\lambda)=\inf_{x\in\mathcal X}L(x,\lambda)$ 为凹函数，所需步骤是最大化它。若 $x^t$ 是精确最小解，$g(x^t)$ 是其超梯度，故
\[
 \lambda^{t+1}=[\lambda^t+\eta_t g(x^t)]_+,
 \qquad\eta_t>0.
\]
约束违反时 $g(x^t)>0$，乘子增加，使后续原变量更新对违反约束付出更大代价。约束 $a^{\T}x\geq b$ 应先写成 $b-a^{\T}x\leq0$，相应更新使用 $b-a^{\T}x^t$。只有写清拉格朗日函数，才可确定更新符号。

\section*{题~\ref{ex:sca}：上界与内近似}
凸函数满足
$h(x)\geq h(x^t)+\nabla h(x^t)^{\T}(x-x^t)$，因此
\[
 g(x)-h(x)\leq
 g(x)-h(x^t)-\nabla h(x^t)^{\T}(x-x^t)
 =:\widetilde f(x;x^t).
\]
应线性化被减去的 $h$。最小化这个凸上界，若它在 $x^t$ 处取等，且新点使上界不增，就有 $f(x^{t+1})\leq f(x^t)$。把 $\widetilde f(x;x^t)\leq0$ 用作约束给出内近似：满足新约束必满足原约束，并保留当前可行点。

单调有界只保证目标值收敛，尚不能保证迭代点收敛。要得到聚点为停点，还需梯度一致、适当的紧致性或有界性、约束正则性以及对子问题的求解精度等条件。停点是一阶条件；非凸目标的不同停点可能有不同目标值，不能由单调性或停点条件推出全局最优。

\section*{题~\ref{ex:wmmse}：先优化接收器，再优化权重}
由不相关性，误差为
\[
 e(u)=|u|^2(|a|^2+I)-2\Rea(u^*a)+1.
\]
配方得到 $u_{\mathrm{opt}}=a/(|a|^2+I)$，$e_{\min}=I/(|a|^2+I)=1/(1+|a|^2/I)>0$。对固定 $e>0$，函数 $\ln v-ve+1$ 关于 $v>0$ 严格凹，其导数为 $1/v-e$，故最大值在 $v=1/e$，值为 $-\ln e$。再最小化误差即得题中恒等式，联合最优权重为 $v_{\mathrm{opt}}=1+|a|^2/I$。

该恒等式给出一种精确的目标重写。多用户中对波束、接收器和权重轮流优化，一般只解联合非凸问题的各个子问题；恒等式并未赋予这些交替迭代全局最优保证。

\section*{题~\ref{ex:ee}：唯一根与能效参数}
分母始终为正，故导数的符号由
\[
 H(p)=\frac{p_c+p}{1+p}-\ln(1+p)
\]
决定。它满足 $H(0)=p_c>0$、$H'(p)=-(p_c+p)/(1+p)^2<0$，且 $H(p)\to-\infty$。因此存在唯一正根 $p_0$，目标先增后减；有上界时最优功率为 $\min(P,p_0)$。内部最优条件为
\[
 (1+p_0)\ln(1+p_0)=p_c+p_0.
\]
当 $p_c=1$ 时 $p_0=\mathrm e-1\approx1.71828<4$，最优能效为 $1/\mathrm e$，单位为归一化的 $\mathrm{nat/J}$。

Dinkelbach 子问题为 $\max_{0\leq p\leq P}\{\ln(1+p)-q(p_c+p)\}$。当 $q>0$，最优解是 $p(q)=\min\{P,\max\{0,1/q-1\}\}$；当 $q=0$ 时取 $P$。每轮更新 $q\leftarrow\ln(1+p(q))/(p_c+p(q))$。以子问题最优值接近零作为停止依据；它为零时恰好满足分式问题的全局最优条件。

\section*{题~\ref{ex:assignment}：指派多面体与耦合约束}
两个整数匹配分别是对角线匹配，收益 $4$，以及交叉匹配，收益 $6$，故后者最优。标准指派问题的约束矩阵是全幺模的，右端为整数；其可行多面体的极点为整数，线性目标因此存在整数最优解。最优解存在不意味着每个最优解都必须是整数。

增加耦合约束 $x_{12}+x_{21}\leq1$。所有满足行列和为 $1$ 的二阶矩阵均可写为
\[
 X=\begin{bmatrix}t&1-t\\1-t&t\end{bmatrix},\qquad 0\leq t\leq1.
\]
新约束为 $t\geq1/2$，目标为 $6-2t$。松弛最优解取 $t=1/2$，收益 $5$；整数可行解只能取 $t=1$，收益 $4$，出现严格间隙。新增约束可能破坏原来的整数极点结构。

\section*{题~\ref{ex:chance}：高斯分位数与失败概率预算}
标量 $a^{\T}x$ 均值为 $\mu^{\T}x$，方差为 $x^{\T}\Sigma x$。取 $\Sigma=LL^{\T}$，并令 $\Phi$ 为标准正态分布函数。所求等价约束为
\[
 \mu^{\T}x+\Phi^{-1}(1-\epsilon)\|L^{\T}x\|_2\leq b.
\]
因 $\epsilon<1/2$，分位数为正，约束为凸二阶锥形式。零方差时随机量退化成常数，原约束直接等价于 $\mu^{\T}x\leq b$，上式仍成立。

对第 $j$ 个事件分配 $0<\epsilon_j<1/2$，令 $\sum_j\epsilon_j\leq\epsilon$，再逐个保证失败概率不超过 $\epsilon_j$。并集界给出联合失败概率不超过总预算，无需事件独立。若两个成功事件独立且概率均为 $0.95$，联合概率仅为 $0.95^2=0.9025$，所以单独的 $0.95$ 保证不足以推出联合的 $0.95$。

\section*{题~\ref{ex:admm}：两个残差都要小}
共识形式常用的残差是 $r^t=x^t-z^t$ 和 $s^t=\rho(z^t-z^{t-1})$，其中 $\rho>0$ 是固定罚参数。原始残差衡量共识约束，对偶残差反映对偶最优性条件中的不一致。当 $r^t=0$ 但 $s^t$ 仍大时，尚不能按标准残差准则停止；阈值应同时含绝对项和相对项，以适应变量维数和尺度。

标准两块凸 ADMM 的典型论证要求 $f,g$ 为闭、真、凸函数，问题存在鞍点，采用相应形式的精确子问题更新并固定正罚参数；原始、对偶残差收敛并不自动保证所有原变量序列唯一收敛。引入非凸神经网络后，凸性和精确求解等条件可能不成立，须针对新模型重证有界性、下降或其他控制关系、聚点的可行性与停点性质。不能直接沿用凸 ADMM 的全局最优保证。
